\section{Signed nonlinear propagation through the interior}
\label{energy:section}

We propagate the signed entry data of Proposition~\ref{prep:signed-entry}
with the amplitude required for the holonomy experiment. The individual
errors need not remain smaller than the linear packet. The useful
cancellation occurs in the difference of those errors. This distinction
requires both a sufficiently sharp energy rate and its preservation on
every intermediate time interval.

Fix the background, branch, and prepared optical coordinates of
Section~\ref{bg:section}. Write
\[
 W_t=L_0+L_1=\partial_t+b^A\partial_A,
 \qquad W_x=L_1-L_0=\partial_x-b^A\partial_A,
 \qquad Z_A=\sqrt a\,\partial_A.
\]
The normalized slots are \(\hat e_i=a^{-1/2}L_i\) and
\(\hat e_A=\partial_A\) in the fixed angular atlas. A tensor and its
array of components in these slots will be distinguished when necessary;
all norms below apply to these component functions. The notation \(D\)
always denotes ordinary derivatives in the fixed optical and angular
coordinates. Thus the commutation order does not add tensor slots or
change their null weights. On a time slice, \(C^r_{\mathrm{mix}}\) and
\(H^r_{\mathrm{mix}}\) include all mixed spacetime derivatives through
order \(r\), with the latter taken in \(L^2(dx\,d\operatorname{vol}_\gamma)\).

Use the fixed margins \(u_2>u_*\), \(c_2>0\), and sufficiently small
\(c_0>0\) of Equation~\eqref{bg:packet-domain}. The retained evolution
has strict time and side inequalities, while its entry face is included:
put
\begin{equation}\label{energy:domain}
 \begin{split}
 L=L_T&=\frac{T+u_*}{2}+c_0,\\
 \Omega_T&=\left\{0\le t<L,\quad
 u+\frac{c_0t}{T}<u_2,\quad
 v+\frac{c_0t}{T}<T+c_2\right\}.
 \end{split}
\end{equation}
The slice intervals in \(x\) have a positive lower length bound, and
\(p_T\) lies within fixed smaller margins for large \(T\). Initial
construction on a slightly enlarged domain is understood. Set
\[
 A=A_T=e^{-\alpha T},\qquad \alpha=\frac{7\kappa}{8},
 \qquad \lambda=\lambda_T=e^{\zeta T}.
\]
The relative array of the linear packet is denoted by \(q_T\).

\begin{proposition}[Signed interior experiment]
\label{energy:signed-experiment}
Fix \(0<\delta<\kappa/32\) and the prescribed finite data seminorm.
Use Proposition~\ref{bg:packet} and Proposition~\ref{prep:signed-entry}
with sufficiently high fixed orders, sufficiently small losses and late
support parameter, and sufficiently accurate linear residuals. For
sufficiently small fixed \(\zeta>0\) and sufficiently large \(T\), the
two entry data have smooth vacuum evolutions \(g_T^\pm\) on the strict
domain \eqref{energy:domain}, in wave-map gauge with target \(g\).
Let \(H_T^\pm\) be the relative arrays of \(g_T^\pm-g\), and put
\[
 F_T^\pm=H_T^\pm\mp q_T,\qquad
 F_T^{\mathrm{odd}}=F_T^+-F_T^-.
\]
Then
\begin{equation}\label{energy:conclusion}
 \begin{aligned}
 \lVert F_T^\pm\rVert_{C^{30}_{\mathrm{mix}}(\Omega_T)}
   &\le e^{(\kappa-2\alpha+\delta)T},\\
 \lVert F_T^{\mathrm{odd}}\rVert_{C^2_{\mathrm{mix}}(\Omega_T)}
   &\le e^{(\kappa-3\alpha+\delta)T}
            +e^{(2\kappa-4\alpha+\delta)T}.
 \end{aligned}
\end{equation}
The exact changes and all three error fields have support in
\(v\ge(1-\nu_0)T-o(1)\). The time slices remain spacelike and the two
moving side faces remain strict future exits. Moreover, for some
\(c>0\),
\begin{equation}\label{energy:small-relative}
 \lVert H_T^\pm\rVert_{C^{30}_{\mathrm{mix}}}
       \le e^{-(\kappa/2+c)T},\qquad
 e^{o(T)}\left(\lVert F_T^{\mathrm{odd}}\rVert_{C^2_{\mathrm{mix}}}
          +\sum_\pm\lVert H_T^\pm\rVert_{C^2_{\mathrm{mix}}}^2\right)
       =o(A_T).
\end{equation}
Here every subexponential factor is at a fixed finite order. No higher
seminorm of the fixed smooth background is required to be small.
\end{proposition}

The coefficient input is precisely Lemma~\ref{bg:coefficients}, obtained
from the uniform low estimates, pure longitudinal jets, and expanding
region estimates in \cite[\CompanionRef{foundation}{in:low}, \CompanionRef{foundation}{in:pure-jets}, \CompanionRef{foundation}{in:higher}, \CompanionRef{foundation}{in:continuation}, \CompanionRef{foundation}{pk2:coefficients}]{Foundation}. We prove the nonlinear transfer below.

The first estimate in Equation~\eqref{energy:small-relative} will let us
choose a nearby observing geodesic and control its boosted frame. The
second will separate the two curvature observations: the odd error and
the quadratic metric error are both smaller than the packet amplitude.
An individual error need not have that property.

Here is the quantitative reason for the proof's organization. We will
construct a nonnegative energy rate \(r_T\), common to both signs and
the background linear equation, with
\[
 \mathcal R_T(t)=\int_0^t r_T(s)\,ds,
 \qquad \mathcal R_T(L)\le(\kappa+\eta)T,
\]
where \(\eta>0\) is arbitrarily small. The individual error starts at
quadratic size and grows by \(e^{\mathcal R_T(t)}\), up to allocated
small losses. Subtracting the signed equations cancels the pure-packet
quadratic term. Besides the pure cubic packet term, the remaining
products contain a packet and an error, or two errors. Since
$\mathcal R_T\ge0$, their combined bound has sizes
\(A^3e^{\mathcal R_T(t)}\) and \(A^4e^{2\mathcal R_T(t)}\).
The propagator from an intermediate time \(s\) to \(t\) must retain
\(e^{\mathcal R_T(t)-\mathcal R_T(s)}\); this preserves those one
and two growth factors after integration. A terminal bound applied
anew to every source would lose the gain.

We first isolate the constant tensor-slot weights in the reduced
operator. They yield the limiting square-root energy rate \(2\kappa\) over
an interval of length \(T/2+O(1)\), independently of the ordinary
commutation order. Lower-order couplings then contribute only a
subexponential factor. After closing the individual evolution, we
estimate the odd forcing with the zeroth-order terms of higher commuted
energies. These control its ordinary derivatives without another
division by \(\sqrt a\).

\subsection{Relative arrays and the exact principal operator}

Let \(H\) be a small symmetric array, and let \(h[H]\) be the
covariant tensor with that array in the normalized slots. Define
\(g_H=g+h[H]\), and write
\[
 G_{ab}=g(\hat e_a,\hat e_b)
       =\begin{pmatrix}0&-1&0\\-1&0&0\\0&0&\gamma_{AB}\end{pmatrix},
 \qquad Q(H)=G+H.
\]
Both \(G\) and its inverse are bounded. Set
\(Z_a=\sqrt a\,\hat e_a=J_a{}^\mu\partial_\mu\); explicitly,
\[
 J_0{}^\mu\partial_\mu=L_0,\qquad
 J_1{}^\mu\partial_\mu=L_1,\qquad
 J_A{}^\mu\partial_\mu=\sqrt a\,\partial_A.
\]
The exact identity
\begin{equation}\label{energy:inverse-relative}
 a g_H^{\mu\nu}=Q(H)^{ab}J_a{}^\mu J_b{}^\nu
\end{equation}
shows which coefficients may be differentiated without creating an
inverse lapse. For example, a coefficient of two angular derivatives
contains \(a\), and a mixed radial--angular coefficient in this frame
contains \(\sqrt a\). Coordinate shift contributions are contained in
the bounded \(J_0\). Smooth inversion of \(Q(H)\) introduces no
negative power of \(a\).

For clarity, this assertion concerns the local hyperbolic reduction.
With \(\nabla=\nabla^g\), the connection difference is exactly
\begin{equation}\label{energy:connection-difference}
 C^\lambda{}_{\mu\nu}
 =\frac12 g_H^{\lambda\rho}
       (\nabla_\mu h[H]_{\nu\rho}
        +\nabla_\nu h[H]_{\mu\rho}
        -\nabla_\rho h[H]_{\mu\nu}).
\end{equation}
If \(V_a\) is the background connection matrix along \(Z_a\), then
the normalized array of \(\sqrt a\,\nabla_{\hat e_a}h[H]\) is
\[
 \mathcal D_aH=Z_aH-V_a^{\mathsf T}H-HV_a.
\]
Thus the normalized connection difference along \(Z_a\) is a smooth
contraction of \(Q(H)^{-1}\) and \(\mathcal D H\). A second covariant
derivative, multiplied by \(a\), is a combination of
\(Z_aZ_bH\), \(Z_aH\), and \(H\); differentiating the factor
\(a^{-1/2}\) produces \(Z_a\log a\), not an inverse lapse.
The differentiated connection matrices and the curvature appear with
the scaled factors specified in Lemma~\ref{bg:coefficients}. Their
uncontracted fixed jets, when needed in nonlinear terms, have
subexponential bounds. In particular,
\[
 D^Ia=a\,e^{o(T)},\qquad
 D^I\sqrt a=\sqrt a\,e^{o(T)},\qquad
 D^I\partial_vb=a\,e^{o(T)}
\]
at every fixed order. These identities also apply to the brackets and
frame coefficients appearing in the preceding calculation.

The Ricci difference formula obtained from
\eqref{energy:connection-difference}, after subtracting the wave-map
gauge derivative, consists of the scalar wave principal term, products
of first covariant derivatives, and algebraic background-curvature
contractions. Normalize its overall constant so that its principal
term is \(a g_H^{\mu\nu}D_\mu D_\nu H\), and denote the resulting
array by \(\mathfrak R(H)\). This is an equivalent reduced equation.
Its linearization at zero is denoted by \(\mathcal L\). It follows
from the displayed identities that
\begin{equation}\label{energy:reduced-structure}
 \mathfrak R(H)_n=(\mathcal L H)_n
       +B^{\mu\nu}(H)D_\mu D_\nu H_n
       +\mathcal Q_n(H,DH),\qquad
 B^{\mu\nu}(H)=a(g_H^{\mu\nu}-g^{\mu\nu}).
\end{equation}
Here \(B(0)=0\), and \(\mathcal Q\) vanishes to second order at
\((H,DH)=0\). More explicitly, it is a finite sum of contractions of
smooth functions of \(H\) with background arrays and field factors
whose total ordinary derivative count is at most two and which contain
at least two field factors. The smooth undifferentiated-field factors
can equivalently be treated by their integral Taylor formula. All
fixed derivatives of these coefficient functions are bounded by
\(e^{o(T)}\) on a fixed small ball of relative arrays. Background
trace reversal, if used to identify the linear operator, is
covariantly parallel and does not change this scalar principal part.

The packet bounds and \eqref{energy:reduced-structure} give, for every
fixed comparison order \(j\),
\begin{equation}\label{energy:packet-residual}
 \lVert\mathfrak R(\pm q_T)\rVert_{H^j_{\mathrm{mix}}}
 \le e^{o(T)}(A\lambda^{-M_1}+A^2\lambda^{P_j})
 \le A^2e^{\eta T}.
\end{equation}
The fixed integer \(P_j\) depends on the base derivative order, not
on the subsequently chosen formal accuracy. We will also require
\(\lVert\mathcal Lq_T\rVert_{H^{60}_{\mathrm{mix}}}
\le e^{-4\kappa T}\). Polynomial slice-volume factors are included
in \(e^{o(T)}\). The number \(\eta>0\) can be prescribed as small
as necessary by the order of choices at the end of the proof.

For one sign, abbreviate \(q=\pm q_T\), \(F=F_T^\pm\), and
\(g'=g_{q+F}\). Subtraction gives the exact equation
\begin{equation}\label{energy:exact-principal}
 \begin{split}
 \mathcal LF+B^{\mu\nu}(q+F)D_\mu D_\nu F
 ={}&-\mathfrak R(q)
  -\bigl(B^{\mu\nu}(q+F)-B^{\mu\nu}(q)\bigr)D_\mu D_\nu q\\
 &-\bigl(\mathcal Q(q+F,Dq+DF)-\mathcal Q(q,Dq)\bigr).
 \end{split}
\end{equation}
After \(j\) ordinary differentiations we retain the exact principal
operator \(a g'^{\mu\nu}D_\mu D_\nu\) on each \(D^IF_n\).
Every other perturbative product contains an error factor and at least
one further packet or error factor, has total field-derivative count
at most \(j+2\), and has no error factor with more than \(j+1\)
derivatives. Indeed, the unique placement with \(j+2\) derivatives
on an error and no differentiated coefficient is precisely the
retained principal term. Every principal commutator has a coefficient
hit. The other terms in \eqref{energy:exact-principal} have either a
second derivative on the known packet or at most first derivatives on
the displayed error factors before commutation. This argument also
covers derivatives of the smooth coefficient functions by the chain
rule.

\subsection{Constant slot weights and positive fluxes}

For a covariant two-tensor entry \(n\), let \(n_i\) count its radial
\(i\)-slots and set
\[
 c_n=\frac\kappa2(n_0-n_1),\qquad |c_n|\le\kappa.
\]
The limiting vector connection matrices in Lemma~\ref{bg:coefficients}
give covariant slot actions \(+c_n\) along \(L_0\) and \(-c_n\)
along \(L_1\). Since \(L_0x=-1/2\) and \(L_1x=1/2\),
\begin{equation}\label{energy:factor-two}
 e^{2c_nx}L_0e^{-2c_nx}=L_0+c_n,
 \qquad e^{2c_nx}L_1e^{-2c_nx}=L_1-c_n.
\end{equation}
Consequently the constant radial tensor operator is
\[
 -(L_0+c_n)(L_1-c_n)-(L_1-c_n)(L_0+c_n),
\]
the radial part of
\begin{equation}\label{energy:conjugation}
 \mathcal C_{n,g'}f
       :=e^{2c_nx}a\square_{g'}^{\mathrm{scal}}(e^{-2c_nx}f).
\end{equation}
The factor two in \eqref{energy:factor-two} is fixed by the definitions
of \(t,x\).

For the background metric, the difference between \(\mathcal L\)
and these diagonal operators is a matrix of first- and zeroth-order
terms on \(L_if_m,Z_Af_m,f_m\). Its coefficients are bounded and
tend to zero in \(\min(t,v,-u)\to\infty\). This assertion includes
the contracted differentiated connection terms and
\(a\operatorname{Riem}(g)\); only the constant diagonal connection
products have been retained. It is exactly the coefficient conclusion
of Lemma~\ref{bg:coefficients}. Replacing \(g\) by \(g'\) in
\eqref{energy:conjugation} and passing to scalar divergence form add
perturbative first- and zeroth-order terms of the type already counted.
For example,
\[
 \sqrt{|\det g'|}=\sqrt{|\det g|}\,
           \left|\frac{\det Q(q+F)}{\det G}\right|^{1/2},
\]
so its new logarithmic derivatives depend smoothly on relative arrays
and their first derivatives. Together with
\eqref{energy:inverse-relative}, this verifies that divergence form
does not introduce an extra inverse lapse.

Fix a small positive mass parameter \(\varepsilon_{\mathrm m}\). For an entry
\(f=D^IF_n\), put
\[
 \widetilde f=e^{-2c_nx}f,\qquad
 \mu^2=\frac{\varepsilon_{\mathrm m}^2}{2a},\qquad
 X_n=e^{4c_nx}W_t.
\]
Use the massive scalar stress tensor of \(\widetilde f\) for \(g'\)
and multiplier \(X_n\). At \(g'=g\), its positive time-slice flux
is exactly
\begin{equation}\label{energy:flux}
 \frac12\int\left\{
 |W_tf|^2+|(W_x-2c_n)f|^2
     +2a\gamma^{AB}\partial_Af\,\partial_Bf
     +\varepsilon_{\mathrm m}^2|f|^2\right\}
                      dx\,d\operatorname{vol}_\gamma.
\end{equation}
Indeed \(W_t\) is normal to the slice and
\(g(W_t,W_t)=-2a\); \(W_x\) is its orthogonal radial tangent,
with squared length \(2a\). Also \(W_tx=Z_Ax=0\) and
\(W_xx=1\). Substitution in the stress tensor gives
\eqref{energy:flux}, including the conjugated radial derivative.
For a small relative change the exact flux is uniformly equivalent
to \eqref{energy:flux}. This is a finite-dimensional comparison of
\(g'/a\) in the basis \(W_t,W_x,Z_A\), whose coefficient matrix
is a small change of a uniformly nondegenerate matrix. It proves
simultaneously that the slices are spacelike and \(W_t\) is future
timelike.

Localize in a finite angular atlas, using cutoffs whose squared sum is
bounded below. Let \(\mathcal E_j\) be the sum of the exact fluxes
at commutation order exactly \(j\), over entries, multiindices, and
charts, and let \(E_j=\sqrt{\mathcal E_j}\). Lower orders will be
included explicitly when needed. The mass terms imply
\begin{equation}\label{energy:mass-sobolev}
 \lVert F(t)\rVert_{C^r_{\mathrm{mix}}}
       \le C_{r,\varepsilon_{\mathrm m}}\sum_{j=0}^{r+2}E_j(t).
\end{equation}
To see this also for temporal derivatives, fix any mixed derivative
\(D^IF\) of order at most \(r\). Its three-dimensional spatial
\(H^2\) norm is controlled by the mass terms of the mixed energies
through \(r+2\). The embedding \(H^2\subset L^\infty\) applies
on each angular chart and radial interval. Uniformity in the interval
length follows by using overlapping intervals of a fixed length and,
at endpoints, the fixed positive minimum slice length. Low equivalence
of \(\gamma\) to a fixed angular metric makes the constants uniform.
No angular lapse conversion occurs in \eqref{energy:mass-sobolev}.

Use individual commutation levels \(0\le j\le60\), with bootstrap
\begin{equation}\label{energy:bootstrap}
 E_j(t)\le e^{-5\kappa T/8}.
\end{equation}
The entry proposition, taken through full mixed order 61, starts this
bootstrap with \(E_j(0)\le A^2e^{\eta T}\). Its pointwise estimates
give the same slice estimates after the polynomial volume loss.
If initial normal jets are instead recovered recursively, the
noncharacteristic reduced equation at \(t=0\) gives precisely these
jets from sufficiently high Cauchy data: the entry lapse is uniformly
nondegenerate. All requisite orders are fixed before \(\zeta\).

The bootstrap and \eqref{energy:mass-sobolev} give small relative
arrays and first derivatives, including the packet. Cone comparison
in \(W_t,W_x,Z_A\) therefore gives, on future causal vectors
parametrized by \(t\),
\begin{equation}\label{energy:cone}
 \frac{du}{dt},\ \frac{dv}{dt}
             \ge-Ce^{-5\kappa T/8+o(T)}.
\end{equation}
For completeness, the background causal inequality in this basis is
\(-2(Vt)^2+2(Vx)^2+\gamma(V_\theta,V_\theta)\le0\), with
\(V_\theta\) the scaled screen component. A small relative metric
change preserves bounded ratios of all components to \(Vt>0\) and
gives \(|Vx|/Vt\le1+C|H|\), which is
\eqref{energy:cone}. Thus each tilted side function in
\eqref{energy:domain} has strictly positive future derivative for
large \(T\). There is no incoming boundary datum on those faces.
Finite propagation, starting from exact early agreement, yields
\begin{equation}\label{energy:known-support}
 \operatorname{supp}H_T^\pm\cup\operatorname{supp}F_T^\pm
       \subset\{v\ge(1-\nu_0)T-\omega_T\},\qquad
 \omega_T=Ce^{-5\kappa T/8+o(T)}T=o(1).
\end{equation}
The packet is supported in a fixed-width neighborhood of \(v=T\).

\subsection{All commutator placements}

We separate background commutators from perturbative products. The
second-coordinate part of the scaled background scalar wave is
\begin{equation}\label{energy:background-principal}
 -2\partial_u\partial_v-2b^A\partial_A\partial_v
                 +a\gamma^{AB}\partial_A\partial_B.
\end{equation}
The \(\partial_u\partial_v\) coefficient is constant. For
\(|I|=j\), a term with exactly one derivative on the shift has form
\[
 (Db^A)L_1\bigl(\partial_AD^{I'}F_n\bigr),
                       \qquad |I'|=j-1.
\]
The parenthesized field has ordinary order \(j\), and its remaining
derivative is radial. It is controlled by \(E_j\) with coefficient
\(C_j|Db|\), with no angular loss. If \(\ell\ge2\) derivatives
hit the shift, the same parenthesized field has order
\(j-\ell+1\le j-1\), so this term belongs to a strictly lower
energy level and its coefficient may cost \(e^{o(T)}\).

For the angular principal term, a single coefficient derivative has
the exact factorization
\[
 D(a\gamma^{AB})
   =a\bigl((D\log a)\gamma^{AB}+D\gamma^{AB}\bigr).
\]
Put one angular derivative inside the order-\(j\) commuted field
and write the other as \(a^{-1/2}Z_A\). The resulting coefficient
is \(\sqrt a\) times a bounded low array. It tends to zero in the
bulk even though \(D\log a\) need not tend to zero. With \(\ell\ge2\)
coefficient hits, the field is at level \(j-\ell+1<j\) and the
coefficient is \(\sqrt a\,e^{o(T)}\).

Write every background first-order term in the directions
\(L_0,L_1,Z_A\). A hit on its displayed coefficient leaves an
energy derivative of a field of order at most \(j-1\); a hit on a
zeroth-order coefficient leaves a lower mass norm. A hit on the
\(b\) inside \(L_0\) can also leave
\((Db)\partial_AD^{j-1}F\), an order-\(j\) mass placement with
small bulk coefficient. Derivatives of the \(\sqrt a\) inside
\(Z_A\) retain that factor and lower the remaining commuted field.
The same reasoning covers angular localizations and chart changes:
\([b^A\partial_A\partial_v,\chi]\) has a \(b\) factor and a
radial derivative, whereas \([a\gamma^{AB}\partial_A\partial_B,\chi]\)
has an \(a\) factor. Differentiated transition matrices act only on
angular slots and preserve \(c_n\). Terms involving the radial
bracket retain \(\partial_vb=a\,O(1)\), so their scaled angular
coefficient is \(a^{-1/2}\partial_vb=O(\sqrt a)\).
Consequently all same-level background remainders have bounded
coefficients, arbitrarily small in the bulk, while coefficients of
strictly lower levels have subexponential bounds.

For perturbative products a rougher estimate suffices, but it must be
used only once. From \eqref{energy:flux}, the bounded shift and the
mass term,
\begin{equation}\label{energy:one-angular-loss}
 \lVert D^{j+1}F\rVert_{L^2}
   \le C_{j,\varepsilon_{\mathrm m}}e^{o(T)}
             \sup_{\Omega_T}a^{-1/2}\sum_{\ell\le j}E_\ell
   \le e^{\kappa T/2+o(T)}\sum_{\ell\le j}E_\ell.
\end{equation}
Indeed \(\partial_A=a^{-1/2}Z_A\), while \(\partial_u\) and
\(\partial_v\) are bounded combinations of \(L_i\) and angular
derivatives. The subtraction of \(2c_n f\) in the radial energy is
paid for by its fixed mass constant. Finally
\(a^{-1/2}\le C e^{\kappa L}=e^{\kappa T/2+O(1)}\).

Place an error factor with the largest error derivative count in
\(L^2\) in each product following \eqref{energy:exact-principal}.
If there is another error factor, every such remaining factor has
order at most \((j+2)/2\le31\). Its supremum is controlled by the
mass energies through order 33, by \eqref{energy:mass-sobolev}, and
costs \(Ce^{-5\kappa T/8}\). Undifferentiated additional error
factors satisfy the same bound. If instead the additional factors
are packets, their derivatives are bounded pointwise by
\(A\lambda^P e^{o(T)}\) for a fixed \(P\), even when a packet has
more derivatives than the selected error. Thus the sole possible
conversion \eqref{energy:one-angular-loss} gives the coefficient
\begin{equation}\label{energy:nonlinear-small-rate}
 e^{\kappa T/2+o(T)}
       \bigl(e^{-5\kappa T/8}+A\lambda^P\bigr)
  =e^{-\kappa T/8+o(T)}
       +e^{(-3\kappa/8+P\zeta)T+o(T)}=o(1).
\end{equation}
This bounds the perturbative source by a small coefficient times
\(\sum_{\ell\le j}E_\ell\). There are never two angular conversions:
the other differentiated errors were estimated by their higher mass
norms, and the retained principal term was the only error placement
with order \(j+2\). The deformation of the actual metric in the
stress identity requires only first relative jets; the scaled
connection calculation above gives the same small coefficients.

We also record the source pairing, since it rules out a different
spurious exponential factor. If \(\mathcal C_{n,g'}f=S_n\), then
the source in the equation for \(\widetilde f\) is
\(a^{-1}e^{-2c_nx}S_n\). Since \(W_tx=0\), its pairing with
\(X_n\widetilde f\) is
\[
 (a^{-1}e^{-2c_nx}S_n)
       (e^{4c_nx}e^{-2c_nx}W_tf)=a^{-1}S_nW_tf.
\]
The spacetime volume element in \((t,x,\theta)\) is
\(2a\sqrt{\det\gamma}\,(1+O(|H|))\,dt\,dx\,d\theta\).
It cancels the remaining \(a^{-1}\). Hence the contribution to the
squared energy is at most
\begin{equation}\label{energy:source-pairing}
 C\lVert S_n\rVert_{L^2}E_j.
\end{equation}
An off-diagonal source involving entry \(m\) is first estimated in
its original variables \(f_m,L_if_m,Z_Af_m\), with its own flux.
The displayed cancellation occurs entirely in entry \(n\); no ratio
\(e^{2(c_n-c_m)x}\) occurs.

\subsection{A common rate on every intermediate interval}

The stress tensor and its divergence, for any metric \(g'\), are
\[
 \begin{split}
 T_{\mu\nu}[\widetilde f]
  &=\partial_\mu\widetilde f\,\partial_\nu\widetilde f
    -\frac12g'_{\mu\nu}
       \bigl(g'^{-1}(d\widetilde f,d\widetilde f)
                    +\mu^2\widetilde f^2\bigr),\\
 \nabla_{g'}^\mu T_{\mu\nu}
  &=(\square_{g'}\widetilde f-\mu^2\widetilde f)
                              \partial_\nu\widetilde f
       -\frac12\partial_\nu\mu^2\,\widetilde f^2.
 \end{split}
\]
Integrate its current with multiplier \(X_n\) between two time
slices. The strict exiting side faces have nonnegative outward energy
flux and may be discarded. To determine the coefficient of the energy
itself, first use a fixed angular metric \(\gamma(\theta)\), set
\(b=0\), and take \(a=a_0e^{-2\kappa t}\). The angular coordinate
coefficients of \(\gamma\) need not be constant. The conjugated
scalar equation is
\[
 \frac12\bigl(-f_{tt}+(\partial_x-2c_n)^2f\bigr)
                +a\Delta_\gamma f=S_n.
\]
Set \(P=f_t\) and \(V=(\partial_x-2c_n)f\). Differentiating
\eqref{energy:flux} and integrating the radial and angular second
derivatives gives
\begin{equation}\label{energy:model-identity}
 \frac{d}{dt}\mathcal E[f]
  =\int\left(-4c_nPV
       -2\kappa a|d_\theta f|_{\gamma^{-1}}^2
       +\varepsilon_{\mathrm m}^2fP-2S_nP\right)
                              dx\,d\operatorname{vol}_\gamma
\end{equation}
up to the favorable boundary flux. For the radial cross term, the two
expressions \(P(\partial_x-2c_n)V\) and
\(V(\partial_x-2c_n)P\), after integration by parts, leave
\(-4c_nPV\). Differentiating the angular coefficient gives
\(a_t|d_\theta f|^2=-2\kappa a|d_\theta f|^2\), and the mass
term gives \(\varepsilon_{\mathrm m}^2fP\). This also verifies the cancellation
between the variable-mass derivative and the mass part of the metric
deformation in the stress formulation.

The angular term is nonpositive, while
\[
 4|c_nPV|\le2|c_n|(|P|^2+|V|^2),\qquad
 \varepsilon_{\mathrm m}^2|fP|
       \le\frac{\varepsilon_{\mathrm m}}{2}
                         (\varepsilon_{\mathrm m}^2|f|^2+|P|^2).
\]
Thus the nonsource rate for the squared energy is at most
\(4\kappa+\varepsilon_{\mathrm m}\), and that for its square root is at most
\(2\kappa+\varepsilon_{\mathrm m}/2\). There is no commutation-order factor
in this constant, since \(c_n\) is the weight of the original two
slots for every \(D^IF_n\).

We justify taking this model as the bulk limit without assuming that
\(\gamma\) converges to a prescribed angular metric. Put
\(Y=(W_t,W_x,Z_A)\) and \(\widehat Q_{ab}=g'(Y_a,Y_b)/a\).
The deformation tensor, derivative of \(e^{4c_nx}\), and derivative
of \(\mu^2\), after dividing their contractions by the energy
density, are smooth algebraic functions of \(\widehat Q\), its
first relative derivatives, the brackets of \(Y\), and the bounded
background low arrays. The required background limits are
\begin{equation}\label{energy:deformation-limits}
 \begin{gathered}
 W_t\log a=-2\kappa+o(1),\qquad W_t\gamma=o(1),\\
 [W_t,W_x]=-2(\partial_vb)^A\partial_A=o(1)Z_A,\qquad
 [W_t,Z_A]=-\kappa Z_A+o(1)Z_B.
 \end{gathered}
\end{equation}
For the third identity use \(\partial_vb=a\,O(1)\); for the last
use \(W_t\log a\) and \(\partial_Ab=o(1)\). The radial metric
block therefore has limiting deformation \(-2\kappa\) times
itself, the angular block has zero limiting deformation, and mixed
deformations tend to zero in relative units. Angular spatial
derivatives remain in divergence form and require boundedness only.
These statements give exactly \eqref{energy:model-identity} in the
limit. The changes in \(\widehat Q\) and its relevant first
derivatives are exponentially small by the bootstrap. Fixing
\(\varepsilon_{\mathrm m}\) first permits every remaining bulk term, including
the bounded first- and zeroth-order linear source coefficients, to be
made arbitrarily small relative to the positive flux. The background
linear equation enjoys the same estimates.

Here is an explicit common rate construction. Fix a desired small
\(\eta>0\). Choose \(\varepsilon_{\mathrm m}\), then a bulk threshold
\(B_0\), so that all same-level square-root rates at the finitely
many used orders are at most \(2\kappa+\eta_0\) on
\(\min(t,v,-u)\ge B_0\), where \(\eta_0\) is as small as needed.
Their rates elsewhere are bounded by a constant \(C_0\). These
constants apply simultaneously to both actual metrics and the
background operator; enlarge \(T\) to absorb
\eqref{energy:nonlinear-small-rate}. On the support
\eqref{energy:known-support}, a point outside the bulk can occur only
at times in
\[
 J_T=[0,B_0]\ \cup\
 \left[\frac{(1-\nu_0)T-\omega_T-B_0}{2},L\right].
\]
Indeed \(v\ge B_0\) for large \(T\), and \(-u=v-2t\).
In particular \(|J_T|\le\nu_0T/2+O(B_0+1)\). Choose \(\nu_0\)
after \(C_0\), sufficiently small, and set for example
\begin{equation}\label{energy:common-rate}
 r_T(t)=2\kappa+\eta_0+C_0\mathbf1_{J_T}(t),\qquad
 \mathcal R_T(t)=\int_0^t r_T(s)\,ds.
\end{equation}
Increasing \(C_0\) if necessary and making the preceding choices
strict gives
\begin{equation}\label{energy:rate-integral}
 \mathcal R_T(L)\le(\kappa+\eta)T.
\end{equation}
It is a single nonnegative pointwise majorant for both individual
equations and the background linear equation at every used order,
not a terminal estimate chosen separately at each level.

Combining the source and commutator bounds yields, in integrated
form or almost everywhere after regularization,
\begin{equation}\label{energy:hierarchy}
 E_j'\le r_T E_j+\sum_{\ell<j}a_{j\ell}(t)E_\ell+S_j(t),
 \qquad a_{j\ell}\ge0,\qquad
 \int_0^L a_{j\ell}(t)\,dt\le e^{o(T)}.
\end{equation}
At a zero of an energy this follows by first using
\((\mathcal E_j+\epsilon^2)^{1/2}\) and then letting
\(\epsilon\downarrow0\). Same-level couplings among entries and
multiindices were summed within \(\mathcal E_j\). The remaining
matrix in \eqref{energy:hierarchy} strictly lowers the ordinary
commutation order.

\begin{lemma}[Finite Volterra propagator]\label{energy:volterra}
For levels \(0,\ldots,K\) satisfying \eqref{energy:hierarchy}, with
fixed \(K\), there is a subexponential \(P_T\), uniform for all
\(0\le s\le t\le L\), such that
\begin{equation}\label{energy:propagator}
 \sum_{j\le K}E_j(t)
 \le P_T\left\{
 e^{\mathcal R_T(t)-\mathcal R_T(s)}\sum_{j\le K}E_j(s)
 +\int_s^t e^{\mathcal R_T(t)-\mathcal R_T(\tau)}
                       \sum_{j\le K}S_j(\tau)\,d\tau\right\}.
\end{equation}
\end{lemma}
\begin{proof}
Let \(\mathbf A(t)=(a_{j\ell}(t))\), with zero diagonal and entries
zero unless \(\ell<j\). Conjugate the inequalities by
\(e^{-\mathcal R_T(t)}\) and use comparison with the nonnegative
linear system. Its time-ordered propagator is the finite sum
\[
 \mathbf U(t,s)=\operatorname{Id}
 +\sum_{m=1}^{K}
   \int_{s<\tau_m<\cdots<\tau_1<t}
       \mathbf A(\tau_1)\cdots\mathbf A(\tau_m)
                             \,d\tau_m\cdots d\tau_1.
\]
Every product of \(K+1\) such matrices is zero: a nonzero entry
would require \(K+1\) strict descents among \(K+1\) levels. This
does not require that the matrices at different times commute. If
\(B_T=\int_0^L\lVert\mathbf A(\tau)\rVert\,d\tau\), then
\(\lVert\mathbf U(t,s)\rVert\le\sum_{m=0}^K B_T^m/m!\).
The right side is \(e^{o(T)}\), since \(K\) is fixed and
\(B_T=e^{o(T)}\); for any desired exponential tolerance use that
tolerance divided by \(K+1\) in each coefficient bound. Variation
of constants and removal of the common conjugation give
\eqref{energy:propagator}. The bound used the whole-interval
\(B_T\), and is therefore uniform in both endpoints.
\end{proof}

\subsection{Individual closure and vacuum continuation}

Apply Lemma~\ref{energy:volterra} with \(K=60\), the entry bounds,
and \(S_j\le A^2e^{\eta T}\) from
\eqref{energy:packet-residual}. By reducing preliminary losses, and
absorbing the finite triangular factor and \(L=O(T)\), we obtain
\begin{equation}\label{energy:individual-time}
 \sum_{j\le60}E_j(t)
       \le A^2 e^{\mathcal R_T(t)+\eta_1T},
               \qquad 0\le t\le L,
\end{equation}
where \(\eta_1>0\) can be made arbitrarily small along with the
tolerance in \eqref{energy:rate-integral}. This is an estimate at
every time. Its terminal exponent is
\(-2\alpha+\kappa+\eta+\eta_1
=-3\kappa/4+\eta+\eta_1\), strictly smaller than
\(-5\kappa/8\) when \(\eta+\eta_1<\kappa/8\). It improves
\eqref{energy:bootstrap}, and \eqref{energy:mass-sobolev} gives the
individual \(C^{30}_{\mathrm{mix}}\) estimate after allocating the
losses within \(\delta\).

For each fixed \(T\), the lapse has a positive minimum on the
finite closed comparison region. The bounded mixed energy norms
therefore give the ordinary local Cauchy norms needed for continuation;
their equivalence constants need not be uniform in \(T\). The
temporal principal coefficient stays nonzero, the relative metric
stays Lorentzian, and the strict side margins persist by
\eqref{energy:cone}. Local restarts are made on the enlarged slice
intervals and restricted by domain of dependence. One can extend
the reduced-equation slice fields smoothly across an artificial side
for this purpose; no constraints on that auxiliary extension are
needed outside the retained dependence domain. A finite maximal
continuation time would contradict these bounds and the local
noncharacteristic existence theorem. Higher smooth orders persist
by differentiating on this fixed regular slab, where the known low
norms control the principal coefficients and each higher finite norm
satisfies a finite linear energy bound. This is the local continuation
argument used in \cite[\CompanionRef{foundation}{sec:interior} and \CompanionRef{foundation}{sec:packet}]{Foundation}.

Finally the exact reduced solutions are vacuum. In the convention
\(\operatorname{Ric}(g')+\nabla^{g'}_{(\mu}\Upsilon_{\nu)}=0\),
the contracted Bianchi identity gives a homogeneous linear wave
equation for the gauge covector \(\Upsilon\). Its initial value
vanishes by the exact entry gauge. Its tangential derivatives then
vanish. In an orthonormal frame with unit normal \(N\), the mixed
normal Einstein constraints give
\(\nabla_N\Upsilon_A=0\), and the normal-normal constraint gives
\(\nabla_N\Upsilon_N=0\): substituting the reduced equation in
\(\operatorname{Ein}(g')(N,\cdot)=0\) leaves these components
with nonzero factors. Thus both gauge Cauchy data vanish. Homogeneous
wave uniqueness on the outflow dependence domain proves
\(\Upsilon=0\) and \(\operatorname{Ric}(g')=0\).

\subsection{Odd forcing with higher mass buffers}

Both individual solutions now exist with
\eqref{energy:individual-time}. Their difference is a known field,
so its support is already given by \eqref{energy:known-support}.
Writing its equation with the background operator does not change
that support.

Set \(\mathcal N(H)=\mathfrak R(H)-\mathcal LH\). At the level of
the finite jet variables in \eqref{energy:reduced-structure}, Taylor's
formula gives
\[
 \mathcal N(H)=\mathcal B_2(H,H)+\mathcal N_{\ge3}(H),
\]
where \(\mathcal B_2\) is a bilinear differential expression of
total derivative count at most two, and the integral remainder has
at least three field factors with the same total derivative count.
The coefficient functions have the previously established fixed-jet
bounds. Subtracting the two exact equations gives
\begin{equation}\label{energy:odd-equation}
 \mathcal LF_T^{\mathrm{odd}}
  =-2\mathcal Lq_T
     -\mathcal N(q_T+F_T^+)+\mathcal N(-q_T+F_T^-).
\end{equation}
The pure packet quadratic terms cancel because
\(\mathcal B_2(q_T,q_T)=\mathcal B_2(-q_T,-q_T)\).
All remaining quadratic terms contain either a packet and an error
or two errors. In the cubic remainder the pure packet term has size
\(A^3\), and every other term contains an error and is bounded by
one of these two quadratic types times a small field factor. This
uses smallness of the total relative fields, not an assumption that
an individual error is smaller than the packet.

Estimate \eqref{energy:odd-equation} through mixed order 10. The
highest derivative in its source has order 12. It is controlled in
unweighted \(L^2\) by the mass term at individual level 12, rather
than by the degenerate angular derivative term at level 11. All
needed supremum factors can, even without optimizing their placement,
be controlled by individual mass energies through level 14. These
levels lie strictly below 60. The ordinary product rule, applied on
each slice to the finite jet expression, therefore gives
\begin{equation}\label{energy:odd-source}
 \lVert\mathcal LF_T^{\mathrm{odd}}(t)\rVert_{H^{10}_{\mathrm{mix}}}
 \le e^{\eta_2T+o(T)}
       \left(A^3e^{\mathcal R_T(t)}
                      +A^4e^{2\mathcal R_T(t)}\right).
\end{equation}
Here \(\eta_2\) is a fixed finite sum of the preliminary losses
and fixed positive frequency powers times \(\zeta\), hence can be
made arbitrarily small. Specifically the packet--error products use
\(A\lambda^P\cdot A^2e^{\mathcal R_T(t)+\eta_1T}\), the
error--error products use
\(A^4e^{2\mathcal R_T(t)+2\eta_1T}\), and
\(A^3\lambda^P\le A^3\lambda^P e^{\mathcal R_T(t)}\).
The negligible \(2\mathcal Lq_T\) is absorbed because its exponent
\(-4\kappa\) is strictly below \(-3\alpha\). Higher Taylor
factors are absorbed by the already proved relative smallness and
the same product estimates. No use of
\eqref{energy:one-angular-loss} occurs in this argument.

The initial odd energy through level 10 is at most
\(A^3e^{\eta T}\) by Proposition~\ref{prep:signed-entry}. Apply
the background version of \eqref{energy:hierarchy} with the same
\(r_T\) and use Lemma~\ref{energy:volterra}. This is legitimate on
the known support of the field in \eqref{energy:odd-equation}; the
background side fluxes also have the favorable sign. Monotonicity of
\(\mathcal R_T\) gives, for every \(t\le L\),
\begin{align}
 \int_0^t e^{\mathcal R_T(t)-\mathcal R_T(s)}
                 A^3e^{\mathcal R_T(s)}\,ds
       &=tA^3e^{\mathcal R_T(t)},\label{energy:duhamel-one}\\
 \int_0^t e^{\mathcal R_T(t)-\mathcal R_T(s)}
                 A^4e^{2\mathcal R_T(s)}\,ds
       &\le tA^4e^{2\mathcal R_T(t)}.\label{energy:duhamel-two}
\end{align}
Thus the odd norm has only one or two final growth factors,
respectively. Applying \eqref{energy:mass-sobolev} at order two,
and absorbing the fixed triangular and polynomial factors, proves
the second line of \eqref{energy:conclusion}. The preserved
intermediate-time difference
\(\mathcal R_T(t)-\mathcal R_T(s)\) is essential in these two
integrals.

\subsection{Finite choices and strict margins}

We finish by specifying the quantifier order. First fix individual
order 60, odd order 10, the desired order 30 for the metric, the
finite data seminorm, and all entry, exterior, constraint, gauge, and
restriction buffers needed for them. Fix \(\delta<\kappa/32\).
Allocate smaller losses so that their finite sum, including twice
the tolerance in \eqref{energy:rate-integral} for the second odd
term, is below \(\delta\). Choose the mass parameter, the bulk
threshold, and then \(\nu_0\) as in
\eqref{energy:common-rate}; the preparation proposition supplies
early agreement for that \(\nu_0\). Choose its exterior losses
within the same allocation. The finitely many positive frequency
powers are now fixed. Take \(\zeta>0\) sufficiently small for these
powers and for \eqref{energy:nonlinear-small-rate}. Only afterwards
increase the finite formal transport and inverse-accuracy orders
until all linear residual bounds used above hold. At every resulting
fixed order, increase \(T\) to absorb constants and subexponential
factors. No simultaneous estimate over an increasing derivative
order is used.

The three error exponents in \eqref{energy:conclusion} are
\[
 -\frac{3\kappa}{4}+\delta,
 \qquad-\frac{13\kappa}{8}+\delta,
 \qquad-\frac{3\kappa}{2}+\delta.
\]
The individual exponent improves the bootstrap by more than
\(3\kappa/32\). The packet's fixed \(C^{30}\) exponent can be
chosen at most \(-7\kappa/8+\kappa/32\); hence, after absorbing a
fixed sum of constants, both total changes are bounded, for example,
by \(e^{-11\kappa T/16}\). This proves the first assertion of
\eqref{energy:small-relative} with \(c=3\kappa/16\) after a further
increase of \(T\). Dividing the two odd bounds by \(A\) gives
exponents at most \(-3\kappa/4+\delta\) and
\(-5\kappa/8+\delta\). Squaring the total-change bound and dividing
by \(A\) gives exponent at most \(-\kappa/2\). All are strictly
negative, and remain so after multiplication by any fixed-order
\(e^{o(T)}\). This proves \eqref{energy:small-relative} and
Proposition~\ref{energy:signed-experiment}.
